% !TEX program = xelatex
\documentclass[a4paper, 14pt]{extreport}

\usepackage{fontspec}
\setmainfont[Ligatures=TeX]{Times New Roman}

% Поддержка русского языка
\usepackage[utf8]{inputenc}
\usepackage[T2A]{fontenc}
\usepackage[russian]{babel}

% Настройка полей (обычно для диплома)
\usepackage[left=3cm, right=1.5cm, top=2cm, bottom=2cm]{geometry}

% Пакеты для абзацев
\usepackage{indentfirst}
\setlength{\parindent}{1.25cm} % Красная строка

% Межстрочный интервал (обычно 1.5)
\usepackage{setspace}
\onehalfspacing

% Настройка заголовков (делаем их адекватного размера)
\usepackage{titlesec}
\titleformat{\chapter}{\normalfont\Large\bfseries}{\thechapter.}{1em}{}
\titleformat{\section}{\normalfont\large\bfseries}{\thesection.}{1em}{}

\usepackage[parentracker=true, backend=biber, style=gost-numeric, sorting=none]{biblatex}
\addbibresource{references.bib} % Подключаем ваш файл с базой

\usepackage{array}
\usepackage{booktabs}
\usepackage{geometry}

\usepackage{tabularx}
\usepackage{enumitem}
\setlist[itemize]{itemsep=0pt,parsep=0pt,partopsep=0pt,topsep=0pt}
\usepackage{float}
\usepackage{longtable}


%Hyphenation rules
%--------------------------------------
\usepackage{hyphenat}
\hyphenation{ма-те-ма-ти-ка вос-ста-нав-ли-вать}
% Английские переносы
\hyphenation{hy-phen-a-tion ex-am-ple com-pu-ter}
%--------------------------------------

\usepackage[unicode]{hyperref}
\setlength{\parskip}{0pt}

\usepackage{graphicx}
\graphicspath{ {./images/} }

\usepackage{booktabs}

\begin{document}
\begin{titlepage}

% \fontsize{12}{14.4}\selectfont

\begin{center}
    \textbf{
    МИНИСТЕРСТВО ЦИФРОВОГО РАЗВИТИЯ, СВЯЗИ И МАССОВЫХ 
    КОММУНИКАЦИЙ РОССИЙСКОЙ ФЕДЕРАЦИИ
    }
\end{center}

\begin{center}
    Ордена Трудового Красного Знамени федеральное государственное \\
    бюджетное образовательное учреждение высшего образования \\
    «Московский технический университет связи и информатики» \\
    (МТУСИ)
\end{center}

\begin{center}
    Кафедра «Математическая кибернетика и информационные технологии»
\end{center}

\vspace{5em}

\begin{center}
    \textbf{
    Дипломная работа
    }
\end{center}

\begin{center}
   на тему:
\end{center}

\begin{center}
    \textbf{
    Исследование применимости обработки данных с использованием методов машинного обучения во встраиваемых системах по сравнению с серверной обработкой
    }
\end{center}

\vspace{10em}

\raggedleft
\begin{tabular}{lr} 
    Выполнил:       & студент группы БВТ2201 \\
                    &  Алимбеков Р. А. \\
                    \cmidrule(lr){2-2}
    Руководитель:   & Симонов С.Е \\
                    \cmidrule(lr){2-2}
\end{tabular}

\vfill
\begin{center}
    Москва \\
    2026
\end{center}
\end{titlepage}

\setcounter{page}{2}
\renewcommand{\contentsname}{Оглавление}
\tableofcontents

\newpage

\chapter{Глава 1}
\section{Эволюция парадигм вычислений: от облачной обработки к TinyML}

Cloud computing - архитектура вычислений, которая подрузомевает передачу  данных от конечного устройства на удаленный сервер для агрегации и обработки этих данных \cite{Mell2011}. Данный подход позволяет преодолеть ограничения IoT устройств \cite{MohammedSadeeq2021}. И мы можем использовать его для передачи сырых  данных с датчикров (например, аудиопотока) на удаленные серверы для агрегации и обработки сложными моделями машинного обучения. Несмотря на доступ к практически неограниченным вычислительным ресурсам серверов, облачный подход имеет ряд критических недостатков.

Технические ограничения: скорость и приватность:
\begin{itemize}
    \item Задержка и джиттер: Даже при наличии 5G, географическая удаленность дата-центров создает задержки, неприемлемые для критически важных приложений. централизованное облако не способно обеспечить детерминированное время отклика для IoT систем \cite{Hazra2024}.
    \item Конфиденциальность: постоянная потоковая передача сырых данных (особенно аудиозаписей из жилых или рабочих помещений) на центральный сервер многократно повышает риски утечки чувствительной информации и нарушения приватности. Отказ от передачи сырого потока в пользу локальной обработки является важным шагом для соблюдения принципов защиты конфиденциальности. \cite{Li2019, Srinivas2025}. 
    \item Зависимость от пропускной способности: Постоянный стриминг данных быстро истощает ресурсы сети \cite{Aazam2014}.
\end{itemize}

Помимо технических ограничений присутствуют и экономические риски:
\begin{itemize}
    \item Vendor Lock-in (Зависимость от провайдера): Использование специфических облачных API делает миграцию на другую платформу практически невозможной без полной переработки системы \cite{OparaMartins2014}.
    \item Стоимость пропускной способности (Bandwidth Cost): Переход в облако требует затрат на высокоскоростные каналы связи. Для приложений с большими объемами данных стоимость передачи может превышать экономию на оборудовании \cite{Dillon2010}
\end{itemize}

В ответ на сетевые и инфраструктурные вызовы сформировалась парадигма Edge Computing (граничные вычисления). Фундаментальная концепция данного подхода заключается в переносе вычислительных мощностей и алгоритмов принятия решений от централизованных облачных серверов непосредственно к границе сети — к источникам генерации данных \cite{Hazra2024}. Локальная обработка информации позволяет радикально снизить задержку отклика системы (Latency), минимизировать зависимость от пропускной способности канала связи и повысить общую отказоустойчивость \cite{Satyanarayanan2017}.

Логичным развитием концепции граничных вычислений в сфере интеллектуального Интернета вещей (AIoT) стало направление TinyML (Tiny Machine Learning). Данная дисциплина находится на стыке машинного обучения и встроенных систем, изучая методы развертывания нейронных сетей на аппаратном обеспечении с экстремально ограниченными ресурсами — микроконтроллерах (MCU) с объемом оперативной памяти менее 1 МБ и энергопотреблением порядка милливатт \cite{warden2019tinyml, Ray2022}. 

Жесткие аппаратные ограничения требуют специфического конвейера разработки. Модели машинного обучения (например, для задачи Keyword Spotting) первоначально обучаются на высокопроизводительных серверах с использованием вычислений с плавающей точкой (Float32). Затем, для адаптации под микроконтроллеры, применяются методы компрессии: прунинг (удаление малозначимых нейронных связей) и пост-тренировочное квантование (снижение разрядности весов до Int8) \cite{Ray2022}. 

Инференс таких оптимизированных моделей осуществляется с помощью специализированных легковесных фреймворков, таких как TensorFlow Lite for Microcontrollers (TFLite Micro) \cite{david2021tensorflow}. Использование TFLite Micro на энергоэффективных чипах семейства ESP32 позволяет реализовать непрерывный анализ аудиопотока и классификацию признаков локально \cite{tflite_micro_docs}. Это обеспечивает полную автономность устройства, гарантирует приватность голосовых данных и практически полностью исключает сетевую нагрузку (Network Load), отправляя во внешнюю среду только готовые управляющие сигналы.

В качестве компромисса между облачной и локальной парадигмами в современных распределенных системах все чаще применяется концепция гибридного интеллекта (Hybrid Edge-Cloud Intelligence) или разделенных вычислений (Split Computing) \cite{Srinivas2025}. Этот архитектурный подход позволяет разделить вычислительный конвейер (pipeline) между маломощным конечным устройством и высокопроизводительным сервером, оптимизируя время отклика в реальном времени без потери аналитической точности.

Механика гибридного метода заключается в строгом функциональном разделении: этап предварительной обработки данных (Data Preprocessing) и извлечения признаков (Feature Extraction) выполняется непосредственно на микроконтроллере, тогда как ресурсоемкий алгоритм инференса сложной модели машинного обучения делегируется серверу. В контексте обработки аудиосигнала это означает, что микроконтроллер локально выполняет захват сырого PCM-потока и проводит детерминированные математические операции — например, оконное преобразование Фурье для вычисления логарифмических мел-спектрограмм (Log-Mel Spectrograms). 

Сформированный тензор акустических признаков, имеющий на порядки меньший объем по сравнению с исходным аудиофайлом, передается по сети на сервер (например, в микросервис на базе FastAPI). На стороне сервера этот компактный вектор подается на вход полноразмерной неквантованной нейронной сети (CNN), которая осуществляет финальную классификацию вероятностей.

Такая архитектура дает синергетический эффект. Во-первых, она обеспечивает высокий уровень конфиденциальности (Privacy-Preserving): по сети передается не разборчивая человеческая речь, а абстрактная математическая матрица признаков, что делает перехват аудио потока бессмысленным \cite{Srinivas2025}. Во-вторых, радикально снижается нагрузка на пропускную способность беспроводной сети и энергопотребление передатчика устройства. При этом система сохраняет доступ к серверным вычислительным мощностям, позволяя использовать тяжелые модели формата Float32, что обеспечивает максимальную точность распознавания, недостижимую для полностью автономных Edge-решений.

\section{Обоснование выбора задачи распознавания ключевых слов и механизмы её реализации в реальном времени}

Для объективного сравнения эффективности облачной, локальной и гибридной вычислительных парадигм необходима задача, которая позволяет наглядно измерить ключевые системные метрики: задержку, сетевую нагрузку, потребление памяти конечного устройства и итоговую точность классификации.

\subsection{Обоснование выбора задачи: обработка аудиосигналов и Keyword Spotting}
Анализ аудиосигналов является оптимальным полигоном для подобного исследования по нескольким причинам. В отличие от данных с простых скалярных датчиков (например, температуры или влажности), сырой аудиопоток обладает высокой плотностью и требует значительной пропускной способности для передачи в облако. С другой стороны, аудиоданные не так избыточны, как видеопоток, который зачастую невозможно обрабатывать локально на микроконтроллерах класса ESP32 в режиме реального времени \cite{Swamy2020, Nowroz}.

Обработка звука идеально раскрывает потенциал гибридного подхода, поскольку этап извлечения признаков (переход от временной области аудиосигнала к частотной, например, расчет мел-спектрограмм) является вычислительно детерминированным, но при этом радикально — на порядки — сжимает объем полезных данных без потери семантической информации.

В качестве конкретной прикладной задачи в данной работе выбрано распознавание ключевых слов (Keyword Spotting, KWS). KWS представляет собой задачу бинарной или многоклассовой классификации, цель которой — непрерывный мониторинг акустического фона и детектирование заранее заданных слов-триггеров \cite{López-Espejo}. Эта задача является классической в экосистеме TinyML, что обеспечивает доступ к стандартизированным датасетам (например, Google Speech Commands \cite{speech_commands_dataset_tfds}) и бенчмаркам (например, MLPerf Tiny \cite{mlperf_tiny_v1_3})позволяет корректно оценивать влияние аппаратных и архитектурных ограничений на качество работы нейросетевых моделей.

\subsection{Методы обеспечения работы KWS в реальном времени}
В реальных условиях эксплуатации системы KWS аудиосигнал поступает не в виде заранее нарезанных аудиофайлов фиксированной длины, а в форме непрерывного потока. Это требует внедрения специфических алгоритмических механизмов на этапе инференса для обеспечения надежного распознавания.\cite{López-Espejo}

Поскольку ключевое слово может быть произнесено в любой момент времени, жесткая сегментация входящего аудиопотока на непересекающиеся блоки (например, строго по 1 секунде) неизбежно приведет к ситуациям, когда слово окажется разделенным между двумя соседними буферами. В таком случае нейронная сеть получит на вход лишь часть фонетических признаков и не сможет корректно классифицировать триггер. 

Для решения этой проблемы применяется метод скользящего окна. В оперативной памяти устройства формируется кольцевой буфер (Ring Buffer) фиксированного размера, хранящий акустические признаки за последнюю секунду. Модель машинного обучения запускается не по мере полного обновления буфера, а с заданным шагом (Stride) — например, каждые 100–500 мс или после накопления определенного количества новых фреймов мел-спектрограммы. Такой подход гарантирует, что искомое слово хотя бы в одной из итераций инференса полностью попадет в рецептивное поле нейронной сети.\cite{López-Espejo}

Однако такой подход приводит к тому, что одно и то же ключевое слово может быть обнаружено несколько раз в последовательных окнах. Для предотвращения повторных срабатываний после детектирования вводится так называемый период невосприимчивости (refractory period), в течение которого система временно приостанавливает обработку новых срабатываний.\cite{López-Espejo}

\subsection{Вычислительный конвейер KWS}

Внутри каждого активного окна захвата система выполняет строго определенный конвейер вычислений, который можно разделить на два основных этапа: предварительную обработку сигнала и нейросетевой инференс. Учитывая целевую платформу (микроконтроллеры семейства ESP32), оба этапа требуют глубокой аппаратной оптимизации.

\subsubsection{Предварительная обработка аудиосигнала}
Подача сырого аудиопотока (PCM-данных) напрямую на вход нейронной сети крайне неэффективна, так как временная область сигнала содержит избыточное количество шума и слабо отражает семантические особенности речи\cite{Palaz2015}. Поэтому первым шагом конвейера является перевод сигнала из временной области в частотную с извлечением репрезентативных акустических признаков.

В KWS-системах для представления аудиосигнала часто используются логарифмические мел-спектрограммы (Log-Mel Spectrograms). Их вычисление включает применение кратковременного преобразования Фурье (STFT) к коротким фреймам сигнала, после чего спектр обрабатывается с помощью набора треугольных фильтров (Mel-filterbank), отражающих особенности восприятия частот человеком. Далее применяется логарифмирование амплитуд для уменьшения динамического диапазона.\cite{warden2019tinyml}

В результате одномерный аудиосигнал преобразуется в двумерное представление признаков, используемое в качестве входа для моделей машинного обучения.

\subsubsection{Инференс модели машинного обучения}

Полученная мел-спектрограмма подается на вход модели машинного обучения — чаще всего это сверточная нейронная сеть (CNN), которая способна эффективно выявлять пространственно-временные паттерны в акустических данных. Сеть анализирует тензор и на выходе формирует вектор вероятностей для каждого из заданных классов (ключевых слов). \cite{Li2022}

\section{TinyML-оптимизации}



\section{Постановка задачи исследования}

Теоретический анализ показывает, что не существует универсального решения для систем KWS во встраиваемых устройствах. Если в идеальных условиях Edge-модель может демонстрировать приемлемые результаты, то при деградации канала связи или повышении уровня фонового шума поведение различных архитектур требует эмпирического исследования.

На основе проведенного анализа формулируется цель исследования: сравнительный анализ эффективности трех архитектур (Cloud, Edge, Hybrid) для задачи распознавания голосовых команд управления на базе микроконтроллера ESP32. Для достижения цели необходимо оценить следующие метрики:

\begin{itemize}
    \item Полную задержку системы (Latency RTT).
    \item Точность классификации команд (Accuracy).
    \item Объем сетевого трафика (Network Load).
    \item Отказоустойчивость систем (Reliability).
\end{itemize}

Цель: Оценка применимости и сравнительный анализ Cloud, Edge и Hybrid архитектур обработки аудиоданных для систем KWS на базе ESP32. Для достижения цели необходимо решить следующие задачи:
\begin{itemize}
    \item Разработать программный комплекс для микроконтроллера ESP32, обеспечивающий захват аудио, локальный препроцессинг (FFT/Log-Mel) и поддержку трех режимов передачи данных.
    \item Спроектировать и обучить нейронную сеть (CNN) для распознавания набора ключевых слов, провести её оптимизацию (квантование) для запуска в среде TFLite Micro.
    \item Реализовать серверную часть (на базе FastAPI/Python) для выполнения инференса в Cloud и Hybrid режимах.
    \item Разработать методику тестирования, включающую эмуляцию различных условий сетевого канала (задержки, потери пакетов) и уровней акустического шума.
    \item Провести экспериментальное сравнение архитектур по метрикам: Latency (задержка), Accuracy (точность), Network Load (трафик) и Reliability (надежность).
    
\end{itemize}

Оценка должна проводиться по разработанной методике, включающей сценарии идеальных условий (для фиксации базовых показателей), реальных условий эксплуатации (наложение шума 50–70 дБ для проверки устойчивости препроцессинга) и деградации сетевого канала (внесение задержек 100–500 мс и потери пакетов 1–5% для проверки гипотезы о преимуществе Hybrid-метода в условиях нестабильной связи).
\input{conclusion}

\newpage

% Библиография (отобразится как список литературы с номерами в квадратных скобках)
\nocite{*}
\printbibliography[title={Список литературы}]


\end{document}

